\section{Current Situation}

\subsection{Problem Statement}

Maritime straits, geographically specific and narrow passages where the majority of the world’s trade, energy, and navy transportation is funneled through, are an arena where there exists an overlap between two conflicting objectives: the need for unobstructed freedom of navigation common to all nations on a global scale, and the desire of particular states to exercise control over their national territories or strategic areas that serve their national interests.

The following five case studies in front of this panel, South China Sea, Strait of Hormuz, Bab el-Mandeb Strait, Panama Canal and Suez Canal, all show a unique feature of the same basic challenge:

\begin{enumerate}
  \item \textbf{Contestation over sovereignty and overlapping claims.} In the South China Sea, six parties contest their sovereign rights over the same area, and the 2016 verdict from the Permanent Court of Arbitration dismissing China's "Nine-Dash Line" claim has been refused by the defendant itself, showing that even if there is a legally binding decision, it does not force any party to accept it without an enforcement mechanism.\footnote{Permanent Court of Arbitration. (2016, July 12). The South China Sea Arbitration (The Republic of the Philippines v. The People's Republic of China) (PCA Case No. 2013-19). \url{https://pca-cpa.org/en/cases/7/}} \footnote{Beech, H. (2021, July 12). China still rejects South China Sea ruling, 5 years on. VOA News. \url{https://www.voanews.com/}}
  \item \textbf{Military escalation driven by states.} Strait of Hormuz shows how an inter-state conflict involving two countries can shut off a chokepoint that carries about one-fifth of global oil and LNG commerce and how no Security Council mechanism exists to stop such action or roll it back quickly.\footnote{U.S. Energy Information Administration. (2025). The Strait of Hormuz is the world's most important oil transit chokepoint. \url{https://www.eia.gov/todayinenergy/}}
  \item \textbf{Disruption by non-state actors.} Bab el-Mandeb shows that the threat to chokepoint security is not confined to states alone; the activities of the Houthi movement have forced a massive rerouting of ships from around the Cape of Good Hope, affecting the global economy in spite of the Houthi movement being a non-state actor.\footnote{United Nations Conference on Trade and Development. (2024). Review of maritime transport 2024. United Nations. \url{https://unctad.org/publication/review-maritime-transport-2024}} \footnote{United Nations Security Council. (2024, January 10). Resolution 2722 (2024) (S/RES/2722). \url{https://undocs.org/S/RES/2722(2024)}}
  \item \textbf{Vulnerabilities of Environment and Infrastructure.} The Panama Canal serves as an example of a more recent form of chokepoint risk, where environmental conditions, namely drought-induced low water levels, can be as effective in limiting shipping capabilities as a military blockade would.\footnote{U.S. Energy Information Administration. (2023, October 31). Drought at the Panama Canal delays energy shipments, increasing shipping costs. \url{https://www.eia.gov/todayinenergy/detail.php?id=60842}}
  \item \textbf{Risk Compounding and Cascading Effects.} An example of the effects of a chokepoint in one location causing further problems for another location is the case of the Suez Canal, and how an event such as the Ever Given grounding in 2021 shows the vulnerability of global shipping routes that cannot have alternatives to these chokepoints.\footnote{Suez Canal Authority. (2021). Investigation report on the grounding of the Ever Given in the Suez Canal on 23 March 2021. \url{https://www.suezcanal.gov.eg/}}
\end{enumerate}

At the core of all five cases there lies a structural flaw: international law, namely the UNCLOS regime, sets the principles of transit passage and navigation freedom; however, it fails to provide an effective mechanism of enforcement of the above mentioned principles if some state or non-state actors decide to breach them.\footnote{United Nations. (1982). United Nations Convention on the Law of the Sea. United Nations Treaty Series, Vol. 1833, No. 31363. \url{https://www.un.org/depts/los/convention_agreements/texts/unclos/unclos_e.pdf}} The Security Council has taken action selectively, mostly in relation to the Yemen conflict via resolutions (United Nations Security Council, 2024), but does not have a mechanism that would prevent, react to and defuse the crisis of chokepoints prior to causing global economic damage.

The effects of such vacuum in the system of international law cannot be limited to the regional scope, due to the small number of chokepoints carrying the bulk of global trade and energy flows, the disruption of any of these chokepoints due to military, non-military or environmental reasons may lead to an increase in shipping costs, destabilization of energy markets and economic instability in the countries which have no direct interest in the conflict. Therefore, the work of this committee should include consideration not only of the conflicts at each individual chokepoint but the overall lack of an effective international security regime.

\subsection{Possible Solutions}

Delegates should explore solutions at three different levels: \textbf{prevention} (to decrease the chance that disputes turn into disruptions), \textbf{reaction} (how to deal with the crisis when it happens), and \textbf{resilience} (how to decrease vulnerability of the global economy to any disruption in any of the chokepoints).

\subsubsection{Prevention}

\begin{itemize}
  \item Improvement in the effectiveness of the UNCLOS rulings and international arbitration decisions, maybe by employing the Security Council's mechanisms to bring about punitive measures against those that do not comply with the decision like PCA in 2016.\footnote{Permanent Court of Arbitration. (2016, July 12). The South China Sea Arbitration (The Republic of the Philippines v. The People's Republic of China) (PCA Case No. 2013-19). \url{https://pca-cpa.org/en/cases/7/}}
  \item Restarting suspended diplomatic initiatives such as the long-awaited ASEAN-China Code of Conduct for the South China Sea in the UN framework.\footnote{Association of Southeast Asian Nations. (n.d.). ASEAN-China dialogue relations. \url{https://asean.org/asean-china-dialogue-relations/}}
  \item Implementing confidence-building measures for conflicting sides, including the establishment of hot lines, procedures for reporting, rules of engagement for coast guard ships and warships operating in the contested waters.
  \item Handling the root cause of disruptions by non-state actors, in particular, the war in Yemen which leads to attacks in the Red Sea by the Houthis, and recognizing that the problem of Bab el-Mandeb's safety goes beyond naval issues.\footnote{United Nations Security Council. (2024, January 10). Resolution 2722 (2024) (S/RES/2722). \url{https://undocs.org/S/RES/2722(2024)}}
\end{itemize}

\subsubsection{Response}

\begin{itemize}
  \item Expansion of existing multilateral naval cooperation structures such as the International Maritime Security Construct, Combined Maritime Forces, or EU-led convoying missions to defend commercial vessels against the use of force without involving any single country's unilaterally.\footnote{Combined Maritime Forces. (2025). About CMF. \url{https://combinedmaritimeforces.com/about/}}
  \item Creation of some standing Security Council system or quick consultative procedure to deal specifically with chokepoints-related crises more effectively than the present case-by-case approach.
  \item Clarification of humanitarian and economic grounds for granting exemptions in case of closure or blockade of a strait in any future instances of that kind.
\end{itemize}

\subsubsection{Resilience}

\begin{itemize}
  \item Promoting investments into other transportation infrastructure such as pipeline capacity circumventing the Strait of Hormuz, as well as extra port capacity that would make each single canal redundant.
  \item Climate change adaptation financing of infrastructure facilities including the Panama Canal as an example of the environment-related maritime security threat.\footnote{U.S. Energy Information Administration. (2023, October 31). Drought at the Panama Canal delays energy shipments, increasing shipping costs. \url{https://www.eia.gov/todayinenergy/detail.php?id=60842}}
  \item Development of international insurance schemes that would allow stabilizing the shipping and energy market situation during temporary disruptions without an economic incentive to use the chokepoint access as a weapon.
  \item Higher standards in cybersecurity and infrastructure resilience of canal and strait operators given their increasing vulnerability to cyberattacks.
\end{itemize}

The delegates should take note of which of these solutions would be more appropriately handled via Security Council resolutions, which rely on voluntary multilateral cooperation outside the Council, and which need to address issues within bilateral or regional conflicts that fall outside the purview of the Council's jurisdiction.

\subsection{Past Actions}

The international community has approached the issue of chokepoint security by using various layers of legal instruments, including a general legal framework, chokepoint-specific treaties, Security Council Resolutions, and ad hoc multinational operations. Since none of the legal regimes is specific only to maritime chokepoints, it is important for the delegates to examine each layer individually since the resolution will have to accommodate all of them simultaneously.

\subsubsection{The General Legal Framework: UNCLOS (1982)}

The UN Convention on the Law of the Sea (UNLOS) established in 1982 is the key international treaty dealing with marine choke points.\footnote{United Nations. (n.d.). United Nations Convention on the Law of the Sea: Table of contents. Division for Ocean Affairs and the Law of the Sea. \url{https://www.un.org/depts/los/convention_agreements/convention_overview_convention.htm}}

\textbf{Part III, Articles 37-44 (Transit Passage).} This section defines the rights of transit passage through straits used for international navigation and is a distinct concept and more elaborate than innocent passage. Article 38 grants the right of transit passage to all ships and aircrafts, even military vessels, to navigate through such straits continuously and with reasonable speed, while Article 44 asserts that the right to transit passage shall not be suspended by the adjacent states, as opposed to innocent passage, which may be suspended by coastal states for security reasons within their territorial waters.\footnote{United Nations. (n.d.). United Nations Convention on the Law of the Sea: Table of contents. Division for Ocean Affairs and the Law of the Sea. \url{https://www.un.org/depts/los/convention_agreements/convention_overview_convention.htm}} This is the law under which the Strait of Hormuz and Bab el-Mandeb are legally open to ships of all flags, and the particular provision that Iran's threats to close the Strait of Hormuz would violate.\footnote{Katzman, K. (2023). Iran's threat to the Strait of Hormuz (CRS Report No. RS20871). Congressional Research Service. \url{https://crsreports.congress.gov/}}

\textbf{Article 19/Article 45 (Innocent Passage).} In case transit passage is inapplicable, UNCLOS maintains the earlier and more restrictive right of innocent passage, which is capable of being temporarily suspended for security reasons and not transit passage. Iran has signed but not ratified UNCLOS, and at times it has claimed the right to deny passage on the basis of the innocent-passage regime and not transit passage.\footnote{Kaye, S. (2019). Freedom of navigation in the Indo-Pacific region. Papers in Australian Maritime Affairs, No. 22. \url{https://www.navy.gov.au/}}

\textbf{Part V (Articles 55-75, Exclusive Economic Zones).} This section explains the exclusive economic zone that is 200 nautical miles where the coastal states have the sovereign rights of resources but not of territory, which is at the core of the law dispute in the South China Sea region, where six parties’ EEZ claims overlap.\footnote{United Nations. (n.d.). United Nations Convention on the Law of the Sea: Table of contents. Division for Ocean Affairs and the Law of the Sea. \url{https://www.un.org/depts/los/convention_agreements/convention_overview_convention.htm}}

\textbf{Annex VII (Compulsory Arbitration).} Compulsory arbitration is available in the UNCLOS between parties as the legal ground by which the Philippines sued China without its agreement in 2013.\footnote{Permanent Court of Arbitration. (2016, July 12). The South China Sea Arbitration (The Republic of the Philippines v. The People's Republic of China) (PCA Case No. 2013-19). \url{https://pca-cpa.org/en/cases/7/}}

One crucial challenge that delegates will have to face: UNCLOS does not have its own enforcement mechanism. It depends upon the compliance of states, diplomacy, and actions taken by other entities including the Security Council in order to make its provisions work – a flaw evident from all case studies discussed in this committee.

\subsubsection{Case Studies}

\subsubsection{The South China Sea: The 2016 Permanent Court of Arbitration Ruling}

The Philippines initiated arbitration proceedings in January 2013, following the procedures outlined in UNCLOS Annex VII despite China's refusal to be a party to or acknowledge such proceedings.\footnote{Permanent Court of Arbitration. (2016, July 12). The South China Sea Arbitration (The Republic of the Philippines v. The People's Republic of China) (PCA Case No. 2013-19). \url{https://pca-cpa.org/en/cases/7/}} It was on July 12, 2016, when the arbitration panel made up of five judges, appointed by the Permanent Court of Arbitration in Hague, according to the UNCLOS' compulsory dispute resolution procedures, issued their final ruling. The main points were that:

\begin{itemize}
  \item China's claim to historic rights within the Nine-Dash Line in the waters of the South China Sea was invalid under UNCLOS, since the ruling indicated that "the Convention defines the limits of maritime entitlements in the South China Sea, which may not exceed the limits specified, thus rendering extinguished any historic rights that China might have had after the Convention's ratification”.\footnote{Permanent Court of Arbitration. (2016, July 12). The South China Sea Arbitration (The Republic of the Philippines v. The People's Republic of China) (PCA Case No. 2013-19). \url{https://pca-cpa.org/en/cases/7/}} \footnote{Center for Strategic and International Studies. (n.d.). Arbitration on the South China Sea: A reading guide. Asia Maritime Transparency Initiative. \url{https://amti.csis.org/}} \footnote{Asia Maritime Transparency Initiative. (2016, July 12). Arbitration ruling: What you need to know. Center for Strategic and International Studies. \url{https://amti.csis.org/arbitration/}}
  \item No part of the contentious area of Spratly Islands, even the largest one named Itu Aba, qualifies for an island status in compliance with UNCLOS provisions that enable creation of 200 nautical miles EEZ and continental shelf.\footnote{Permanent Court of Arbitration. (2016, July 12). The South China Sea Arbitration (The Republic of the Philippines v. The People's Republic of China) (PCA Case No. 2013-19). \url{https://pca-cpa.org/en/cases/7/}}
  \item China violated the Philippines’ sovereignty by intruding on its oil exploration, stopping Philippine fishing in its own EEZ and not stopping the fishing activities of Chinese fishing vessels in its own EEZ.\footnote{U.S.-China Economic and Security Review Commission. (2016). 2016 annual report to Congress. \url{https://www.uscc.gov/annual-report/2016-annual-report-congress}}
  \item China failed in its commitment to protect the marine environment in accordance with UNCLOS by undertaking massive land reclamation that did serious damage to coral reefs.\footnote{U.S.-China Economic and Security Review Commission. (2016). 2016 annual report to Congress. \url{https://www.uscc.gov/annual-report/2016-annual-report-congress}}
  \item The Scarborough Shoal has been proven to be a traditional fishing ground available for fishermen from various nations such as the Philippines, China, and Vietnam.\footnote{Republic of the Philippines, Foreign Service Institute. (n.d.). The West Philippine Sea dispute. Department of Foreign Affairs. \url{https://www.fsi.gov.ph/}}
\end{itemize}

The ruling is legally binding on both parties according to Articles 296 and 11 of UNCLOS Annex VII, the Chinese government’s refusal to take part in the process does not impact the validity of the decision but only its practical implementation.\footnote{Permanent Court of Arbitration. (2016, July 12). The South China Sea Arbitration (The Republic of the Philippines v. The People's Republic of China) (PCA Case No. 2013-19). \url{https://pca-cpa.org/en/cases/7/}} China has dismissed the decision as having no binding effect and has engaged in land reclamations, militarizations of artificial islands, and coast guard activities in territories identified as Philippine waters by the court, with frequent incidents occurring at the Second Thomas Shoal, which has been consistently characterized in the context of the 2016 ruling by the Philippine government.\footnote{U.S.-China Economic and Security Review Commission. (2016). 2016 annual report to Congress. \url{https://www.uscc.gov/annual-report/2016-annual-report-congress}} The ruling exemplifies in the best way the legal-validity vs. practical-enforcement dichotomy presented in this guide.

\subsubsection{Bab el-Mandeb and the Red Sea: The Yemen Sanctions and Maritime Security Resolutions}

The Security Council has put together an incremental sanctions regime and maritime security regime in relation to the Yemen crisis for more than a decade, as described in the resolutions of the Security Council itself and the records of its Yemen Sanctions Committee:

\begin{itemize}
  \item \textbf{Resolution 2140 (2014)}, created the first Yemen sanctions regime, which was adopted by consensus, involved targeted people whose activities present a threat to the maintenance of peace and security of Yemen under sanctions for one year in the form of freezing their assets and travel ban, and establishment of the 2140 Sanctions Committee and an experts’ group.\footnote{United Nations Security Council. (2014, February 26). Resolution 2140 (2014) (S/RES/2140). \url{https://undocs.org/S/RES/2140(2014)}} \footnote{United Nations Security Council. (n.d.). Security Council Committee established pursuant to resolution 2140 (2014). \url{https://www.un.org/securitycouncil/sanctions/2140}}
  \item \textbf{Resolution 2216 (2015)} pushed for compliance with resolution 2201 (2015) by all parties of Yemen, but mainly the Houthis, under Chapter VII with 14 votes in favor and abstaining vote of Russia, as well as making it obligatory to cease any unilateral actions, leave areas occupied during the crisis and hand over the military hardware acquired from Yemeni Armed Forces.\footnote{United Nations Security Council. (2015, April 14). Resolution 2216 (2015) (S/RES/2216). \url{https://undocs.org/S/RES/2216(2015)}} The most significant step was the imposition of embargo on arms trade regarding some individuals and organizations which gave the ground for stopping arms transportation via sea to the Houthis.\footnote{United Nations Security Council. (2015, April 14). Resolution 2216 (2015) (S/RES/2216). \url{https://undocs.org/S/RES/2216(2015)}} \footnote{United Nations Security Council. (n.d.). United Nations Security Council consolidated list. \url{https://www.un.org/securitycouncil/content/un-sc-consolidated-list}}
  \item \textbf{Resolution 2624 (2022)} was adopted with 11 in favour to 0 against and four abstentions from Brazil, Ireland, Mexico, and Norway which maintained the measures of asset freeze, travel ban, and targeted arms embargo for the first time in resolutions 2140 and 2216; renewed the mandate of the Panel of Experts; and designated the Houthis as one of the entities subject to the targeted arms embargo.\footnote{United Nations Security Council. (2022, February 28). Resolution 2624 (2022) (S/RES/2624). \url{https://undocs.org/S/RES/2624(2022)}} \footnote{United Nations Security Council. (2022, February 28). Resolution 2624 (2022) (S/RES/2624). \url{https://undocs.org/S/RES/2624(2022)}} The states abstaining from voting had issues regarding the definition of the "terrorist group," as well as humanitarian repercussions due to such designation.\footnote{United Nations. (2022, February 28). Security Council renews Yemen sanctions regime, designating Houthis as subject to arms embargo (SC/14805). \url{https://press.un.org/en/2022/sc14805.doc.htm}}
  \item \textbf{Resolution 2722 (2024)} came into effect following the increase in attacks by the Houthis on ships passing through the Red Sea. It deplored attacks on civilian and commercial ships, called for their immediate end, took into account the right of member states, under international law, to protect their ships from such attacks, attacks that could threaten the freedom of navigation, and called on member states to develop maritime security capabilities to ensure the protection of the coastal and port states of the Red Sea and Bab el-Mandeb area.\footnote{United Nations Security Council. (2024, January 10). Resolution 2722 (2024) (S/RES/2722). \url{https://undocs.org/S/RES/2722(2024)}} This was the closest legal cover for the multinational maritime operations that ensued.
  \item \textbf{Resolution 2675 (2023)} re-implemented the Yemen sanctions regime established in Resolution 2140.\footnote{United Nations Security Council. (2023, November 9). Resolution 2675 (2023) (S/RES/2675). \url{https://undocs.org/S/RES/2675(2023)}}
\end{itemize}

It is important that the delegates understand the legal framework that has been utilized in this case; instead of adopting an integrated resolution for maritime security, the council has adopted an approach involving a series of modifications to a sanctions regime for internal conflicts, a reactive approach, not a proactive approach, which relies on periodic extension and not permanent authorization.

Alongside these resolutions, the multinational naval responses that have been initiated include the formation of the International Maritime Security Construct in 2019 to protect shipping in the region and subsequent naval escort and interception operations carried out in response to the Houthi threat post 2023 in the absence of any Chapter VII mandate for naval escort missions.

\subsubsection{The Panama Canal: The Torrijos-Carter Treaties (1977)}

The Torrijos-Carter Treaties were signed by President of the U.S. Jimmy Carter and General Omar Torrijos of Panama on September 7, 1977. There are two treaties that were submitted for consideration to the United Nations Treaty Series.

\begin{itemize}
  \item \textbf{The Panama Canal Treaty} is known for the repeal of the Hay-Bunau-Varilla Treaty passed in 1903 concerning the rule of the United States over the Panama Canal and Canal Zone until December 31, 1999.\footnote{Panama Canal Treaty, Pan.-U.S., September 7, 1977, United Nations Treaty Series, Vol. 1161, No. 18379 (registered 1980). \url{https://treaties.un.org/}}
  \item \textbf{The Treaty Concerning the Permanent Neutrality and Operation of the Panama Canal ("the Neutrality Treaty")} has no termination date. The first article states that the canal is declared to be a neutral international transit waterway; the second article states that both countries will ensure that the Canal shall be safe and accessible to peaceful transit by vessels of all nations "on terms of entire equality," without any discrimination in transit conditions and charges.\footnote{Treaty Concerning the Permanent Neutrality and Operation of the Panama Canal, Pan.-U.S., September 7, 1977, United Nations Treaty Series, Vol. 1161, No. 18380 (registered 1980). \url{https://treaties.un.org/}} A 1977 Statement of Understanding issued by President Carter and General Torrijos, sent via official diplomatic channels, specifies that both countries can take actions - even through the use of force - to protect neutrality and the peaceable transit, but disclaims any right of the US to intervene in the internal affairs of Panama.\footnote{U.S. Department of State, Office of the Historian. (n.d.). Milestones in the history of U.S. foreign relations: The Panama Canal and the Torrijos-Carter Treaties. \url{https://history.state.gov/milestones/}}
\end{itemize}

The importance of such a structure is immediate, given the continued relevance of the Neutrality Treaty and its ongoing impact on the continued involvement of the United States in maintaining the neutrality of the Canal – thus becoming an essential part of the legal context for assessing U.S. pronouncements regarding the security of the Panama Canal and Chinese business interests in the region.

\subsubsection{The Suez Canal: The Convention of Constantinople (1888)}

The Convention of Constantinople is the oldest legal document discussed here and it was concluded on October 29, 1888 by the then major powers of Europe and the Ottoman Empire. As per Article I of the Convention, "The Suez Maritime Canal shall always be free and open, in time of war as in time of peace, to every vessel of commerce or of war, without distinction of flag, and the exercise of the right of blockade against vessels sailing therein shall at no time be asserted."\footnote{Suez Canal Authority. (n.d.). Convention of Constantinople 1888. \url{https://www.suezcanal.gov.eg/}} The convention has been in force since the nationalization of the canal in 1956.

The fact that Egypt did not revoke the international nature of the Canal after nationalization is demonstrated by its unilateral declaration in 1957, submitted to the UN Secretariat, wherein it stated that it was still its “unaltered policy and firm purpose...to respect the terms and the spirit of the Constantinople Convention of 1888 and the rights and obligations arising therefrom”.\footnote{Sabel, R. (2023). The Suez Canal: Past, present, and future. In C. Lutmar \& Z. Rubinovitz (Eds.), The Suez Canal: Past lessons and future challenges (pp. 21–38). Palgrave Macmillan.} This convention does not have a Security Council monitoring mechanism and does not require periodic re-negotiation.
